% Part: applied-modal-logics
% Chapter: temporal-logic
% Section: extra-operators

\documentclass[../../../include/open-logic-section]{subfiles}

\begin{document}

\olfileid{aml}{tl}{ext}

\olsection{સમયસંબંધી તર્ક માટે વધારાના સંક્રિયકો}

ભૂતકાળ અને ભવિષ્ય માટેના એકસ્થાની સંક્રિયકો ઉપરાંત સમયસંબંધી
તર્કપદ્ધતિઓમાં ક્યારેક ``ત્યારથી'' અને ``ત્યાં સુધી'' દર્શાવતા
દ્વિસ્થાની સંક્રિયકો~$\Since $ અને~$\Until$ પણ હોય છે. એટલે
સમયસંબંધી તર્કની ભાષામાં $\Since $ અને~$\Until$ ઉમેરીએ છીએ અને
સમયસંબંધી !!{formula}ની વ્યાખ્યામાં નીચેની શરત ઉમેરીએ છીએ:

\begin{itemize}
\item[] જો $!A$ અને $!B$ એ !!{formula}s હોય, તો $(\Since  !A !B)$ અને
  $(\Until !A !B)$ બંને !!{formula}s છે.
\end{itemize}

પછી આ સંક્રિયકોનું અર્થવિજ્ઞાન નીચે મુજબ છે:

\begin{defn}\ollabel{defn:since-until}
  નિદર્શ~$\mModel M$માં \emph{!!a{formula}~$!A$નું~$t$એ સત્ય હોવું}:
  \begin{enumerate}
  \item\ollabel{defn:sub:mmodels-since} \indcase{!A}{\Since  !B
    !C}{$\mSat{M}{\indfrm}[t]$ ત્યારે અને માત્ર ત્યારે જ જ્યારે
    $t' \prec t$ ધરાવતા કોઈ $t' \in T$ માટે $\mSat{M}{!B}[t']$ હોય,
    અને $t' \prec s \prec t$ ધરાવતા દરેક $s$ માટે
    $\mSat{M}{!C}[s]$ હોય}
   \item\ollabel{defn:sub:mmodels-until} \indcase{!A}{\Until !B
    !C}{$\mSat{M}{\indfrm}[t]$ ત્યારે અને માત્ર ત્યારે જ જ્યારે
    $t \prec t'$ ધરાવતા કોઈ $t' \in T$ માટે $\mSat{M}{!B}[t']$ હોય,
    અને $t \prec s \prec t'$ ધરાવતા દરેક $s$ માટે
    $\mSat{M}{!C}[s]$ હોય}
  \end{enumerate} 
\end{defn}

\sourcecorrection{OLREA-140}{અહીંની બંને અર્થવ્યાખ્યામાં
મધ્યવર્તી બિંદુઓ માટે કડક અસમાનતાઓ છે. એટલે $!C$નું સત્ય
માત્ર $t'$ અને $t$ વચ્ચેના ખુલ્લા અંતરાલમાં જરૂરી છે;
આ શરતથી $!C$ પોતે $t$એ સત્ય હોવાનું મળતું નથી.
બંને બિંદુઓ વચ્ચે કોઈ બિંદુ ન હોય તો મધ્યવર્તી શરત
આપમેળે સંતોષાય છે. નીચેના સહજ વાંચનનો અર્થ આ ચોક્કસ
અર્થવ્યાખ્યા મુજબ લેવો.}

$\Since  !B !C$નો સહજ અર્થ છે: ``$!B$ સાચું હતું ત્યારથી
$!C$~સાચું રહ્યું છે.'' અને $\Until !B !C$નો સહજ અર્થ છે:
``$!B$ સાચું થશે ત્યાં સુધી $!C$~સાચું રહેશે.''

\end{document}
